\documentclass{researchnote}

\usepackage{lipsum}

% --- shared macros (would live in common/macros.tex) ---
\newcommand{\dd}{\mathrm{d}}

\title{An example note title\\on a second line}
\runningtitle{Example note}
\author{Author Name}
\affiliation{Institution Name}
\date{1 January 2026}
\notetags{tag one \textperiodcentered{} tag two \textperiodcentered{} tag three}
% link back to the origin of the note: git repository, agent session, web page...
\notesource{https://example.org/link/to/source}

\begin{document}

\maketitle

\begin{abstract}
This is the summary box. It holds a short description of what the note is about, two to five
sentences long. The link below it points back to where the note came from.
\end{abstract}

\section{First section}

\lipsum[1][1-5] \citet{Doe2020} is a textual citation and \citep{Doe2020} a parenthetical one;
several at once give \citep{Doe2020,Lee2024,Roe2018}. A numbered equation looks like this:
\begin{equation}
  f(x) = \sum_{k=0}^{n} a_k\, x^k .
  \label{eq:poly}
\end{equation}

\begin{tcdefinition}{A widget}{widget}
A \emph{widget} is any object $w$ with a size $|w| \geq 0$ and a label $\ell(w)$.
\end{tcdefinition}

\section{Second section}

\subsection{A subsection}

\lipsum[2][1-4] Equation~\eqref{eq:poly} can be referred to, as can
Figure~\ref{fig:example}.

\begin{tcproposition}{A first claim}{claim}
For any two widgets $u$ and $v$,
\begin{equation}
  |u \oplus v| \leq |u| + |v| .
  \label{eq:claim}
\end{equation}
\end{tcproposition}

\begin{figure}[t]
  \centering
  \includegraphics[width=0.5\linewidth]{example-image-a}
  \caption{An example figure. The caption describes what is shown in one or two sentences.}
  \label{fig:example}
\end{figure}

\subsection{Another subsection}

\lipsum[3][1-3]

\section{Third section}

\begin{tctheorem}{A main result}{main}
Under Proposition~\ref{prop:claim}, for all real $a$ and $b$,
\begin{equation}
  \bigl((a+b)+1\bigr)^2 = a^2 + 2ab + b^2 + 2a + 2b + 1 .
\end{equation}
\end{tctheorem}

\begin{proof}
Expand step by step:
\[
\begin{WithArrows}[displaystyle]
  \bigl((a+b)+1\bigr)^2
    & = (a+b)^2 + 2(a+b) + 1 \Arrow{we expand} \\
    & = a^2 + 2ab + b^2 + 2a + 2b + 1 . \qedhere
\end{WithArrows}
\]
\end{proof}

A second reference, to a book \citep[see chapter~2 of][]{Roe2018}. A list of points:

\begin{itemize}
  \item \textbf{First point.} \lipsum[4][1-2]
  \item \textbf{Second point.} \lipsum[5][1-2]
\end{itemize}

\begin{marker}
\textbf{Warning box.} Use this for caveats or things that are easy to get wrong.
\end{marker}

\bibliography{example}

\end{document}
